    \setlength{\parindent}{0pt}
    \onehalfspacing
    \justifying
    
\noindent\textbf{GLOSSÁRIO} \textit{(opcional)}

Glossário é uma ferramenta linguística/terminológica para apresentação do significado de palavras ou preferencialmente termos. A palavra tem sua existência no contexto geral, enquanto o termo conta com um significado particular, dependendo da área em que se encontra vinculado. Na apresentação dos termos e significado, sugere-se que sejam apresentadas fontes, contudo, em caso de significados atribuídos pelo próprio autor ou significados contextuais, a apresentação da fonte não é totalmente aplicável.\\
\textit{Exemplos}:

\textit{Bibliografia}: no contexto de trabalhos acadêmicos, a bibliografia muitas vezes está relacionada com o conjunto de insumos que determinada pessoa se apropriou, para a elaboração de uma pesquisa. A bibliografia pode ser lida também como um conjunto ordenado de materiais sobre determinado assunto. Em uma pesquisa, a bibliografia não precisa necessariamente ter sido citada para demonstrar relevância no trabalho, ela muitas vezes é vista também como sugestões para se aprofundar em determinado tópico, diferentemente das referências.

\textit{Dissertação}: no Brasil, a dissertação é comumente relativa à materialização para a defesa de uma pesquisa à nível de mestrado.

\textit{Referência}: a principal diferença da referência para a bibliografia, é a citação. As referências, como seu próprio termo diz, são documentos citados para a composição ou apresentação de determinada fonte. Em um trabalho acadêmico, às referências se configuram com uma relação de materiais citados ao longo de uma pesquisa, são a sua base teórica ou metodológica.